
\documentclass[12pt]{amsart}

%==================== Packages ====================
\usepackage[margin=1in]{geometry}
\usepackage{amsmath,amssymb,amsthm}
\usepackage{mathtools}              % paired delimiters, \coloneqq, etc.
\usepackage{enumitem}               % control list labels
\usepackage{bm}                     % bold math: \bm{v}
\usepackage{xcolor}
\usepackage[colorlinks=true,linkcolor=blue,urlcolor=blue]{hyperref}  % load last

\allowdisplaybreaks                 % let long align blocks break across pages
\setlist[enumerate,1]{label=(\alph*), leftmargin=*}  % problem parts as (a), (b), (c)

%==================== Number systems ====================
\newcommand{\N}{\mathbb{N}}
\newcommand{\Z}{\mathbb{Z}}
\newcommand{\Q}{\mathbb{Q}}
\newcommand{\R}{\mathbb{R}}
\newcommand{\C}{\mathbb{C}}
\newcommand{\F}{\mathbb{F}}

%==================== Linear algebra ====================
\DeclareMathOperator{\Ran}{Ran}         % range (Treil); \ker already exists
\DeclareMathOperator{\rank}{rank}
\DeclareMathOperator{\nullity}{nullity}
\DeclareMathOperator{\spn}{span}        % NOT \span: that's a TeX primitive
\DeclareMathOperator{\tr}{tr}
\DeclareMathOperator{\col}{Col}
\DeclareMathOperator{\row}{Row}
\DeclareMathOperator{\proj}{proj}
\DeclareMathOperator{\id}{id}
\newcommand{\Mat}[2]{M_{#1\times #2}(\F)}   % \Mat{m}{n} -> M_{m x n}(F)
\newcommand{\Pn}[1][n]{\mathcal{P}_{#1}}    % \Pn -> P_n, \Pn[3] -> P_3
\newcommand{\tp}{^{\mathsf{T}}}             % A\tp -> A^T
\newcommand{\inv}{^{-1}}                    % A\inv -> A^{-1}
\newcommand{\mat}[1]{\begin{bmatrix}#1\end{bmatrix}}
\newcommand{\pmat}[1]{\begin{pmatrix}#1\end{pmatrix}}

%==================== Analysis / topology ====================
\DeclareMathOperator{\diam}{diam}
\DeclareMathOperator{\dist}{dist}
\DeclareMathOperator{\interior}{int}    % NOT \int: that's the integral sign
\DeclareMathOperator{\sgn}{sgn}
\DeclareMathOperator{\supp}{supp}
\newcommand{\closure}[1]{\overline{#1}}
\newcommand{\pre}{^{\mathrm{pre}}}          % f\pre(V), Pugh's preimage
\newcommand{\nbhd}[2]{N_{#1}(#2)}           % \nbhd{r}{p} -> N_r(p), Rudin
\newcommand{\limn}{\lim_{n\to\infty}}
\newcommand{\eps}{\varepsilon}
\newcommand{\dd}{\,\mathrm{d}}              % \int f(x)\dd x
\newcommand{\restr}[2]{\left.#1\right|_{#2}}  % \restr{f}{A} -> f|_A
\newcommand{\seq}[2][n]{\left(#2\right)_{#1\in\N}}  % \seq{p_n} -> (p_n)_{n in N}

%==================== Delimiters (mathtools) ====================
% \abs{x} normal size; \abs*{x} auto-sizes; \abs[\big]{x} manual size
\DeclarePairedDelimiter{\abs}{\lvert}{\rvert}
\DeclarePairedDelimiter{\norm}{\lVert}{\rVert}
\DeclarePairedDelimiter{\floor}{\lfloor}{\rfloor}
\DeclarePairedDelimiter{\ceil}{\lceil}{\rceil}
\DeclarePairedDelimiter{\inner}{\langle}{\rangle}
\DeclarePairedDelimiterX{\set}[1]{\lbrace}{\rbrace}{#1}
\DeclarePairedDelimiterX{\setof}[2]{\lbrace}{\rbrace}{#1 \,:\, #2}  % \setof{x\in\R}{x>0}

%==================== Logic shortcuts ====================
\newcommand{\contra}{\Rightarrow\!\Leftarrow}   % contradiction marker
\newcommand{\st}{\text{such that }}

%==================== Theorem environments ====================
\theoremstyle{plain}
\newtheorem*{theorem*}{Theorem}
\newtheorem*{lemma*}{Lemma}
\newtheorem*{claim*}{Claim}
\theoremstyle{definition}
\newtheorem*{definition*}{Definition}
\theoremstyle{remark}
\newtheorem*{remark*}{Remark}

\renewcommand{\qedsymbol}{$\blacksquare$}       % black box

% Problem header using the textbook's exercise number: \begin{problem}{2.5}
\newenvironment{problem}[1]{\par\medskip\noindent\textbf{#1.}\ }{\par\smallskip}
% For computational answers that aren't proofs
\newenvironment{solution}{\begin{proof}[Solution]}{\end{proof}}

%==================== Document ====================
\title{Homework 6}
\author{Antony Perez}
\date{\today}

\begin{document}
\maketitle

\begin{problem}{1}
Let $(X,d)$ be a metric space, let $K \subset X$ be compact, and let $\mathcal{U}$ be an open cover of $K$. Prove that there exists a number $\delta > 0$ \st for every $x \in K$, there is some $U \in \mathcal{U}$ for which 
\[
B(x,\delta)\cap K \subset U.
\] 
Such a number $\delta$ is called a \textit{Lebesgue number} for the cover $\mathcal{U}$ (relative to $K$).
\end{problem}

\begin{proof}
    Suppose for contradiction that there is no Lebesgue number. Let $a_n \in K$. Thus, for any $n$, $B(a_n,1/n)\cap K \not \subseteq U.$ By compactness, there exists a convergent subsequence such that $a_{n_k} \to p$. By closedness, $p \in K$. $\mathcal{U}$ covers $K$, so $p \in U_0$, where $U_0 \in \mathcal{U}$. Since $U_0$ is open, there exists an $r>0$ such that $B(p,r) \subseteq U_0$. Now choose a $k$ large enough so that $d(a_{n_k}, p) < r/2$ and $1/n_k < r/2$. This is true, because $1/n_k \to 0$ as $n_k \to \infty$. Now let $x \in B(a_{n_k},1/n_k)\cap K$. By triangle inequality, $d(x,p) \le d(x,a_{n_k}) + d(a_{n_k},p) < 1/n_k + r/2 < r$. Therefore, $x \in B(p,r)$. $x$ was abitrary, so $B(a_{n_k},1/n_k) \cap K \subseteq U_0$. This yields a contradiction. Thus, $\delta$ exists. 
\end{proof}

\begin{problem}{2}
Let $X$ be a metric space, let $K \subset X$ be a nonempty compact set, and let $p \in X \setminus K $ ($K^c$). Define 
\[
d(p,K) = \inf\{d(p,x):x \in K\}.
\]
Prove that there exists $x_0 \in K$ \st 
\[
d(p,K) = d(p,x_0) > 0.
\]
In particular, show that a point outside a compact set has positive distance from that set. 
\end{problem}

\begin{proof}
    $p \in K^c$. Since $K$ is compact, $K$ is closed and bounded, and its complement is open. Thus, $\exists r > 0$ \st $B(p,r) \cap K = \emptyset$. This proves that for $\forall x\in K,$ $d(p,x) \ge r > 0$. Now we must show that $d(p,K)$ is attainable (infimum exists). Let $f:X \to \R$ be defined by $d(p,x)$, where $x \in X$. Using the reverse triangle inequality, $|d(x,p)-d(y,p)| \le d(x,y)$. Let $d(x,y) < \epsilon$, where $\epsilon > 0$. Thus, by reverse triangle inequality, $f$ is continuous. Now, define a sequence $x_n \in K$ \st $d(p,x_n) \to d(p,x_0)$. We can see that $d(p,x_0) \le d(p,x_n) < d(p,x_n) + 1/n$. A subsequence converges to $x_0 \in K$, where convergence and the limit point containment are given by the compactness. Thus, $d(p,x_0) = \lim d(p,x_{n_k})$. Thus, $x_0$ does exist. 
\end{proof}

\begin{problem}{3}
Let $C_0 = [0,1]$. Obtain $C_1$ by removing the open middle third $(1/3,2/3)$ from $C_0$. Inductively, obtain $C_n$ from $C_{n-1}$ by removing the open middle third of each interval in $C_{n-1}.$ The Cantor set is 
\[
C = \bigcap_{n=0}^{\infty} C_n.
\]
\begin{enumerate}
  \item[(a)]   Prove that $C_n$ is the union of $2^n$ pairwise disjoint closed intervals, each of length $3^{-n}$.
  \item[(b)]   Prove that the Cantor set $C$ is nonempty and compact.
  \item[(c)]   Prove that $C$ contains no nonempty open interval. Hence $C$ has empty interior. 
\end{enumerate}

\end{problem}

\begin{proof}
    Use induction. We see $C_0 = [0,1]$, so we have one closed interval with length 1. Now assume $C_n$ satisfies these conditions, and deduce $C_{n+1}$. To obtain $C_{n+1}$, we take $C_n$ and cut out the middle of each interval. This produces two disjoint intervals for each mother interval, thus giving us $2 \cdot 2^n = 2^{n+1}$, and its length would be $3^{-n}/3 = 3^{-(n+1)}$. Since the intervals in $C_n$ are pairwise disjoint, the resulting intervals of $C_{n+1}$ are also pairwise disjoint. Hence, $C_{n+1}$ is the union of $2^{n+1}$ disjoint intervals, each of length $3^{-(n+1)}$. \\
    Let $E$ be the set of endpoints of all the $C^n$ intervals. An endpoint of $C_n$ remains an endpoint for $C_{n+1}$, and so on. This is because removing the middle third of the intervals never touches the endpoints. So $E \subseteq C$. $E$ is infinite, so $C$ is nonempty and infinite. Each $C_n$ is a finite union of closed bounded intervals, so it is compact by Heine-Borel Theorem. Since $C \neq \emptyset$, and each $C_n$ is closed and bounded in $[0,1]$, so $C$ is compact. \\
    Now suppose for contradiction that there does exist an open interval in $C$. Let $c$ be in this open interval. Thus, for $\epsilon > 0$, $B(c,\epsilon) \subseteq C$. Since $C \subseteq C_n$, $B(c,\epsilon) \subseteq C_n$. Thus, $c$ must lie in one of those intervals. Choose $1/3^n$ such that $1/3^n < 2\epsilon$ (obtained by extracting diameter of ball). At whatever $n$ chosen, $B(c,\epsilon) \not \subseteq C_n$. But $B(c,\epsilon) \subseteq C$, and $C \subseteq C_n$. Thus, this is a contradiction. Therefore, $C$ contains no nonempty open interval. 
\end{proof}

\begin{problem}{4}
Recall that a ternary expansion of $x \in [0,1]$ has the form
\[x = \sum_{n=1}^{\infty} \frac{a_n}{3^n} = 0.a_1a_2a_3\ldots \ a_n \in \{0,1,2\}\]
\begin{enumerate}
    \item[(a)] Prove that $x \in C$ if and only if $x$ has a ternary expansion using only the digits 0 and 2. Be careful about numbers that have two ternary expansions, for example 
    \[\frac{1}{3}=0.1_3 = 0.02222\ldots_3\]
    \item[(b)] Use part (a) to prove that the Cantor set is uncountable.
\end{enumerate}
\end{problem}

\begin{proof}
    Let $x \in C$. 
\end{proof}

\begin{problem}{5} Let $X$ be a metric space.
    \begin{enumerate}
    \item[(a)]
    Let $\{E_\alpha\}_{\alpha \in A}$ be a collection of connected subsets of $X$ such that 
    \[\bigcap_{\alpha \in A} E_\alpha \neq \emptyset.\]
    Prove that $\bigcup_{\alpha \in A} E_\alpha$ is connected. 
    \item[(b)] Let $E \subset X$ be connected and suppose
    \[
    E \subset F \subset \overline{E}.
    \]
    Prove that $F$ is connected. In particular, deduce that the closure of $\overline{E}$ of a connected set is connected.
\end{enumerate}
\end{problem}

\begin{proof}
    Suppose $\bigcup_{\alpha \in A} E_\alpha = A \sqcup A^c,$ where $A$ is a proper clopen set. Since $\bigcap_{\alpha \in A} E_\alpha \neq \emptyset$, let $p \in \bigcap_{\alpha \in A} E_\alpha$. From this, let $p \in A$. For each $E_\alpha$, we see $p \in A \cap E_\alpha$. Since each $E_\alpha$ is connected and nonempty, $A \cap E_\alpha$ must be all of $E_\alpha$. This is because of $E_\alpha$ being clopen, by inheritance principle, forces $E_\alpha$ to be the empty set or all of $E_\alpha$. Hence $E_\alpha \subseteq A$ for every $\alpha$. So $\bigcup_{\alpha \in A} E_\alpha \subseteq A$, and $A^c = \emptyset$, contradicting properness. 
\\
Prove the contrapositive. Suppose $F = A \sqcup A^c$, where $A$ is a proper clopen subset in $F$. Let $K = A \cap F$ and $L = A^C \cap F$. This gives clopen subsets in $S$ by inheritance principle. Connectedness of $E$ forces either $K$ or $L$ to be the empty set. Suppose $K = \emptyset$. Let $p \in A$. Create a ball $B(p,r)$, where $r > 0,$ such that $B(p,r) \cap E = \emptyset$. But $p \in F \subseteq \overline{E}$, and every ball around $p$ intersects $E$. Contradiction. 
\end{proof}

\end{document}
